\section{Annotation}

\subsection{На русском:}

Данная работа посвящена исследованию моделирования рыночных шоков с использованием агентно-ориентированных моделей.
В работе описывается постановка задачи – анализ поведения рыночной системы при возникновении внезапных изменений условий торговли.
Основное внимание уделяется разработке экспериментов с использованием симулятора, реализованного в модуле AgentBasedModel, анализу влияния различных параметров (число трейдеров, рыночных мейкеров, интенсивность событий) на динамику системы и валидации результатов моделирования.
Полученные результаты позволяют выявить закономерности формирования волатильности, изменения ликвидности и адаптивного поведения агентов в условиях рыночных потрясений.
Работа имеет практическую значимость для разработки стратегий стабилизации финансовых рынков и оценки рисков в нестабильных экономических условиях

\subsection{In English:}

This work is dedicated to the study of market shock modeling using agent-based models.
The paper outlines the research problem—analyzing the behavior of a market system under sudden trading condition changes.
The primary focus is on developing experiments using a simulator implemented in the AgentBasedModel module, analyzing the impact of various parameters (such as the number of traders, market makers, and event intensity) on the system dynamics, and validating the simulation results.
The obtained results reveal patterns in volatility formation, liquidity changes, and the adaptive behavior of agents under market stress.
The work is of practical importance for developing strategies to stabilize financial markets and assess risks in unstable economic conditions
